\documentclass[../main]{subfiles}
\begin{document}

\chapter{謝辞}
\label{thankyou}

\thispagestyle{empty}

\newpage

% textlint-disable ja-technical-writing/no-mix-dearu-desumasu

本論文を締めくくるにあたり，ご指導，ご協力を頂きました全ての方々に，この場をお借りし，深く御礼の意を評します．

まず，本研究の指導教員である東京大学大学院 新領域創成科学研究科 人間環境学専攻教授 山下 淳先生の手厚いご指導，ご支援に心より感謝申し上げます．お忙しい中，ミニ論文や学会の予稿の執筆時の大変的確なご指摘を，ご助言をいただきました．また，ミーティングの際には研究の方向性についてご意見をいただき，自身の研究について深く考えるきっかけをいただきました．そして，懇親会や懇親会や学会等の場において，お気にかけていただきました．先生のおかげで充実した研究生活を送ることができました．改めて心より感謝申し上げます．

東京大学大学院 新領域創成科学研究科 人間環境学専攻準教授 安 琪先生には，研究の進め方や発表の仕方，論文やポスターの書き方など基礎から丁寧にご教授いただきました．毎週のミーティングでは，研究の方向性や結果の議論を行い，的確なアドバイスをいただきました．大変お忙しい中で，いつでも親身に相談に乗っていただいたおかげで研究へのモチベーションを持ち続けて研究を進めていくことができました．心より感謝申し上げます．

東京大学大学院 新領域創成科学研究科 人間環境学専攻特任講師 濵田 裕幸先生には，グループミーティングや本研究の実験の際に適切なご助言や温かいご指導をいただきました．顔を合わせるたびにお声をかけていただき，研究に関する相談もさせていただきました．心より感謝申し上げます．

東京大学大学院 新領域創成科学研究科 人間環境学専攻特任研究員 菊地 謙様には，実験時のカメラシステムの構築など実験の設定を考えるところからご助言いただき，実験を行う際も進行にご助力いただき，実験結果の解釈についても1から教えていただきました．心より感謝申し上げます．

東京大学大学院 工学系研究科 附属総合研究機構特任助教 伊賀上 卓也さんには，実験時のカメラ配置や設定，実験解析に必要となるステレオ計測をご教授いただきました．また，研究の方向性や発表を見ていただき率直な意見や的確な修正の提案をしていただきました．心より感謝申し上げます．

千葉工業大学 未来ロボット技術研究センター主席研究員 入江 清先生には，実験準備のご協力や卓球に関する知見から多くのご助言をいただきました．また，解析に用いたステレオ計測に関するシステムの開発にもご協力いただきました．心より感謝申し上げます．

九州大学 工学研究院機械工学部門准教授 西川 鋭先生には，実験準備のご協力やミーティング，学会でお会いした際に様々なご指摘，ご助言をいただきました．また，今後の研究の方向性に関しても相談に乗っていただきました．心より感謝申し上げます．

順天堂大学 保健医療学部教授 松田 雅弘先生，順天堂大学 保健医療学部講師 北原 エリ子先生には，実験準備のご協力や今後の車椅子卓球の研究の方向性についてご意見をいただきました．心より感謝申し上げます．

株式会社オカムラ 七野 一輝様，七野様のコーチ 山岸 護様には，海外遠征でお忙しい中お時間をいただき実験にご協力いただきました．心より感謝申し上げます．

研究室の先輩である Rukiye Aydinさん，三上 恭平さんにはお忙しい中本論文のチェッカーを引き受けていただきました．心より感謝申し上げます．

研究室の先輩である 髙村 悠真さん，若松 宥太さんには，卒業論文の執筆に必要不可欠なTeX 環境の構築を手伝っていただきました．また研究外の部分でもよく気にかけていただき，研究室生活を気軽に過ごすことができました．心より感謝申し上げます．

研究室の先輩である 安藤 大生さん，星井 智仁さんには，研究室で顔を合わせるたびに研究で困っていることはないかなどよく声をかけていただきました．心より感謝申し上げます．

同じ研究グループに所属していた 陳 童さん，川端 陸さん，早瀬 瑞華さん，内山 裕稀さん，徐 晨峻さんには，毎週のグループミーティングや研究会などで多くのアドバイスをいただきました．スライドの構成や研究の方針について丁寧に教えてくださったり，昼食に誘っていただいたりと大変お世話になりました．心より感謝申し上げます．

同じ研究室のスタッフ，先輩方には，研究に向かう姿勢のお手本を見せていただき，僕のモチベーションとなりました．心より御礼申し上げます．

同級生である 白井 柊梧さん，東海 和弥さん，山口 滉平さんには，研究室内や学科のイベントなどで会う機会も多く，お互いの研究の話をすることで刺激を受けることができました．心から感謝申し上げます．修士課程でもよろしくお願いします．

秘書の大野朋美さん，小沼博子さんには，学会や出張，実験に用いた会議室や卓球台を借りるための事務手続をして頂き，研究に専念できました．また財布を落とした際もいろいろ助けていただきました．心より感謝申し上げます．

東大GFの皆様には，普段のサークル活動以外でも様々な場所で時間を共に過ごしました．皆様に支えていただいたおかげでとても充実した1年を送ることができました．本当にありがとうございました．

最後に，これまで私を支え，応援してくれた家族や友人に深く感謝申し上げます．
本当にありがとうございました．

\begin{flushright}
  2026年2月 明石 直也
\end{flushright}

% textlint-enable ja-technical-writing/no-mix-dearu-desumasu

\end{document}
